% Part: incompleteness
% Chapter: representability-in-q
% Section: representable-comp

\documentclass[../../../include/open-logic-section]{subfiles}

\begin{document}

\olfileid[id]{inc}{req}{rpc}
\olsection{Fungsi yang Dapat Direpresentasikan dalam $\Th{Q}$ Dapat Dihitung}

Kita akan membuktikan bahwa setiap fungsi yang dapat direpresentasikan
dalam~$\Th{Q}$ dapat dihitung. Pertama-tama kita harus menetapkan sebuah lema
tentang fungsi-fungsi yang dapat direpresentasikan dalam~$\Th{Q}$.

\begin{lem}\ollabel{lem:rep-q}
  Jika $f(x_0, \dots, x_k)$ dapat direpresentasikan dalam~$\Th{Q}$, terdapat
  !!a{formula}~$!A_f(x_0, \dots, x_k, y)$ sedemikian sehingga
  \[
  \Th{Q} \Proves !A_f(\num{n_0}, \dots, \num{n_k}, \num{m})
  \quad\text{jika dan hanya jika}\quad m = f(n_0, \dots, n_k).
  \]
\end{lem}

\begin{proof}
  Bagian ``jika'' ialah
\olref[int]{defn:representable-fn}\olref[int]{defn:rep:a}. Bagian ``hanya
jika'' terlihat sebagai berikut. Andaikan $\Th{Q} \Proves !A_f(\num{n_0},
\dots, \num{n_k}, \num{m})$, tetapi $m \neq f(n_0, \dots, n_k)$. Misalkan
$l = f(n_0, \dots, n_k)$. Menurut
\olref[int]{defn:representable-fn}\olref[int]{defn:rep:a}, $\Th{Q}
\Proves !A_f(\num{n_0}, \dots, \num{n_k}, \num{l})$. Menurut
\olref[int]{defn:representable-fn}\olref[int]{defn:rep:b}, $\Th{Q}$
membuktikan $\lforall[y][(!A_f(\num{n_0}, \dots, \num{n_k}, y) \lif
\num{l} = y)]$. Dengan menggunakan logika dan asumsi bahwa $\Th{Q} \Proves
!A_f(\num{n_0}, \dots, \num{n_k}, \num{m})$, kita memperoleh bahwa
$\Th{Q} \Proves \eq[\num{l}][\num{m}]$. Di sisi lain, menurut
\olref[bre]{lem:q-proves-neq}, $\Th{Q} \Proves
\eq/[\num{l}][\num{m}]$. Jadi $\Th{Q}$ tidak konsisten. Namun, hal itu
mustahil karena $\Th{Q}$ dipenuhi oleh model standar (lihat
\olref[int][def]{def:standard-model}), $\Sat{N}{\Th{Q}}$, dan teori yang
dapat dipenuhi selalu konsisten menurut Teorema Kesahihan
(\tagrefs{prfAX/{fol:axd:sou:cor:consistency-soundness},
prfSC/{fol:seq:sou:cor:consistency-soundness},
prfND/{fol:ntd:sou:cor:consistency-soundness},
prfTab/{fol:tab:sou:cor:consistency-soundness}}).
\end{proof}

\begin{lem}
Setiap fungsi yang dapat direpresentasikan dalam $\Th{Q}$ dapat dihitung.
\end{lem}

\begin{proof}
Pertama-tama mari kita berikan gagasan intuitif mengapa hal ini benar. Untuk
menghitung~$f$, kita melakukan hal berikut. Daftarkan semua kemungkinan
!!{derivation}s~$\delta$ dalam bahasa aritmetika. Hal ini dapat dilakukan
secara mekanis. Untuk setiap kemungkinan itu, periksa apakah ia merupakan
!!a{derivation} dari !!a{formula} berbentuk~$!A_f(\num{n_0}, \dots,
\num{n_k}, \num{m})$ (!!{formula} yang mewakili $f$ dalam~$\Th{Q}$ dari
\olref{lem:rep-q}). Jika ya, maka $m = f(n_0, \dots, n_k)$ menurut
\olref{lem:rep-q}, dan kita telah menemukan nilai~$f$. Pencarian ini berakhir
karena $\Th{Q} \Proves !A_f(\num{n_0}, \dots, \num{n_k},
\num{f(n_0, \dots, n_k)})$, sehingga pada akhirnya kita menemukan suatu
$\delta$ dari jenis yang tepat.

Uraian ini belum sepenuhnya tepat karena prosedur kita bekerja pada
!!{derivation}s dan !!{formula}s, bukan hanya pada bilangan, dan kita belum
menjelaskan secara persis mengapa ``mendaftarkan semua kemungkinan
!!{derivation}s'' dapat dilakukan secara mekanis. Namun, seperti yang telah
kita lihat, suku, !!{formula}s, dan !!{derivation}s dapat dikodekan dengan
bilangan G\"odel. Kita juga telah memperkenalkan model komputasi yang tepat,
yakni fungsi rekursif umum. Selain itu, telah kita lihat bahwa relasi
$\Prf[\Th{Q}](d,y)$, yang berlaku jika dan hanya jika $d$ adalah bilangan
G\"odel dari !!a{derivation} atas !!{formula} bernomor G\"odel~$y$ dari
aksioma-aksioma~$\Th{Q}$, bersifat rekursif (primitif). Fungsi-fungsi rekursif
primitif lain yang kita perlukan ialah $\fn{num}$
(\olref[art][trm]{prop:num-primrec}) dan $\fn{Subst}$
(\olref[art][sub]{prop:subst-primrec}). Dari fungsi-fungsi ini, $f$ dapat
didefinisikan dengan pencarian tak berbatas; dengan demikian, $f$ bersifat rekursif.

Pertama, definisikan
\begin{multline*}
  A(n_0, \dots, n_k, m) = \\
  \fn{Subst}(\fn{Subst}(\dots\fn{Subst}(\Gn{!A_f}, \fn{num}(n_0), \Gn{x_0}),\\ \dots),
  \fn{num}(n_k),  \Gn{x_k}), \fn{num}(m), \Gn{y})
\end{multline*}
Ini tampak rumit, tetapi sebenarnya hanya fungsi
$A(n_0, \dots, n_k, m) =
\Gn{!A_f(\num{n_0}, \dots, \num{n_k}, \num{m})}$.

Sekarang, perhatikan relasi~$R(n_0, \dots, n_k, s)$ yang berlaku jika
$(s)_0$ adalah bilangan G\"odel dari !!a{derivation} dalam~$\Th{Q}$ yang
membuktikan $!A_f(\num{n_0}, \dots, \num{n_k}, \num{(s)_1})$:
\[
R(n_0, \dots, n_k, s) \quad\text{jika dan hanya jika}\quad
\Prf[\Th{Q}]((s)_0, A(n_0, \dots, n_k, (s)_1))
\]
Jika kita dapat menemukan~$s$ sedemikian sehingga $R(n_0, \dots, n_k, s)$
berlaku, berarti kita telah menemukan pasangan bilangan---$(s)_0$ dan
$(s)_1$---sedemikian sehingga $(s)_0$ adalah bilangan G\"odel dari
!!a{derivation} yang membuktikan
$A_f(\num{n_0}, \dots, \num{n_k}, \num{(s)_1})$.
Jadi, mencari~$s$ sama seperti mencari pasangan $d$ dan $m$ dalam pembuktian
informal. Fungsi yang dapat dihitung dan ``mencari'' $s$ semacam itu dapat
didefinisikan dengan pencarian tak berbatas pada fungsi reguler. Perhatikan
bahwa $R$ bersifat reguler: untuk setiap $n_0$, \dots, $n_k$, terdapat
!!a{derivation}~$\delta$ yang menyaksikan
$\Th{Q} \Proves !A_f(\num{n_0}, \dots, \num{n_k},
\num{f(n_0, \dots, n_k)})$, sehingga
$R(n_0, \dots, n_k, s)$ berlaku bagi
$s = \tuple{\Gn{\delta}, f(n_0, \dots, n_k)}$. Jadi, kita dapat menuliskan
$f$ sebagai
\[
f(n_0,\dots,n_{k}) = (\umin{s}{R(n_0, \dots, n_k, s)})_1.
\]
\end{proof}

\end{document}
